\section*{Methode und Grenzen}
\addcontentsline{toc}{section}{Methode und Grenzen}

\subsection*{Was gepr\"uft ist}

Jede Zahl dieses Dokuments stammt aus einer der im Quellenverzeichnis
genannten Prim\"arquellen und steht genau einmal in \texttt{scripts/kennzahlen.py},
zusammen mit ihrer gedruckten Seite. Abgeleitete Gr\"o{\ss}en werden gerechnet und
nicht getippt, damit Text und Tabellen einander nicht widersprechen k\"onnen.
Vier Skripte erzwingen das: \texttt{pruefe\_seiten.py} pr\"uft f\"ur jeden Rohwert,
dass er auf der zitierten Seite tats\"achlich vorkommt; \texttt{check\_literals.py}
verbietet hartkodierte Zahlen im Flie{\ss}text; die beiden Fu{\ss}notenpr\"ufer
halten den Belegapparat konsistent. Beim ersten Durchgang waren
\num{23} von \PruefungAngaben{} Seitenangaben falsch, ausnahmslos die aus dem
Zusammenhang geschlossenen. Deshalb wird geprüft und nicht vertraut.

\subsection*{Drei Farben, drei Verantwortlichkeiten}

Das Dokument trennt sichtbar, wer f\"ur eine Aussage einsteht.
\textbf{Schwarz} ist belegt: Zahl, Tabelle, Quelle, Seite. \bk{\textbf{Blau}
ist eine Einsch\"atzung im laufenden Text, die der Verfasser pr\"uft und
verantwortet.} \vd{\textbf{Rot} ist das Votum in Teil~\ref{sec:verdikt}:
eine aus den Befunden abgeleitete Empfehlung unter erkl\"arten Annahmen,
ausdr\"ucklich kein Beleg und keine Zusage.} Wer nur die schwarzen Stellen
liest, erh\"alt den gepr\"uften Sachstand ohne jede Wertung.

\subsection*{Was nicht gepr\"uft ist}

Gepr\"uft ist die \emph{Herkunft} der Zahlen, nicht die Richtigkeit der
Schl\"usse. Drei Dinge leistet dieses Dokument ausdr\"ucklich nicht:

\begin{enumerate}
\item \textbf{Es prognostiziert nicht.} Die Fortschreibung in Teil~\ref{sec:flow}
ist eine Kapazit\"atsrechnung unter erkl\"arten Annahmen, keine Vorhersage. Die
entscheidende Gr\"o{\ss}e, das k\"unftige organische Wachstum, wird nicht
gesch\"atzt, sondern in drei vom Unternehmen selbst berichteten Auspr\"agungen
durchgerechnet.
\item \textbf{Es gewichtet die Szenarien nicht.} Welches der drei eintritt, ist
eine Einsch\"atzung \"uber Preiszyklen und Integrationserfolg. Die
Prim\"arquellen tragen sie nicht, und das Dokument nimmt sie nicht vorweg.
\item \textbf{Es ersetzt keine Anlageentscheidung.} Anlagehorizont,
Risikotragf\"ahigkeit und Portfoliozusammenhang stehen in keiner
Prim\"arquelle und kommen hier nicht vor. Das Votum in
Teil~\ref{sec:verdikt} nennt einen Schwellenkurs und zwei beobachtbare
Ausl\"oser; ob und in welcher Gr\"o{\ss}e gehandelt wird, entscheidet der
Verfasser.
\end{enumerate}

\subsection*{Wann diese Analyse veraltet}

Sie ruht auf dem Jahresbericht 2025 und dem Zwischenbericht zum
30.\,Juni 2026. Die tragende Gr\"o{\ss}e, das organische Umsatzwachstum, wird
quartalsweise ver\"offentlicht. Mit jedem neuen Zwischenbericht ist zuerst
Abschnitt~\ref{sec:organisch} fortzuschreiben und danach zu pr\"ufen, ob das
Urteil in Teil~\ref{sec:urteil} noch tr\"agt.
